\documentclass[11pt]{article}

\usepackage[margin=1in]{geometry}
\usepackage{amsmath,amssymb,amsthm}
\usepackage{hyperref}

\newtheorem{theorem}{Theorem}
\newtheorem{lemma}{Lemma}
\newtheorem*{remark}{Remark}

\begin{document}

\title{Holomorphic local equations for symplectic simple normal crossings divisors}
\author{}
\date{}
\maketitle

\begin{abstract}
We prove that a symplectic simple normal crossings divisor admits, after the
allowed cohomologous deformation of the symplectic form, a tame almost complex
structure for which the divisor is locally cut out by $J$-holomorphic coordinate
functions.
\end{abstract}

\section*{Conventions}

A \emph{symplectic simple normal crossings divisor} in a symplectic manifold
$(M,\omega)$ is a finite union
\[
D=\bigcup_{i\in S}D_i
\]
of closed embedded symplectic hypersurfaces such that, for every
$I\subset S$, the stratum
\[
D_I=\bigcap_{i\in I}D_i
\]
is either empty or a transverse symplectic submanifold of real codimension
$2|I|$, with the usual positive normal-crossings orientation.  This positivity
is necessary for any holomorphic normal-crossings model.

We write $z:U\to \mathbb C$ is $J$-holomorphic if
\[
dz\circ J=i\,dz .
\]

\section*{Two standard results}

We use two standard theorems from the theory of symplectic normal crossings
divisors.

\begin{theorem}[Orthogonalization; McLean--Farajzadeh Tehrani--Zinger]
\label{thm:orthogonalization}
Let $(M,\omega)$ be closed symplectic and let $D\subset M$ be a symplectic
simple normal crossings divisor.  Then there is a smooth path of symplectic
forms $\omega_t$, $t\in[0,1]$, such that
\[
\omega_0=\omega,\qquad [\omega_t]=[\omega]\in H^2(M;\mathbb R),
\]
each stratum of $D$ remains symplectic for $\omega_t$, and, for $\omega_1$, the
branches of $D$ are pairwise symplectically orthogonal along every stratum.

This is the orthogonalization theorem for symplectic normal crossings divisors;
see McLean \cite[Section 5]{McLean}, especially the orthogonalization lemmas
there, and the regularization formulation of Farajzadeh Tehrani--Zinger
\cite[Theorem 1.2]{FTZ}.  The deformation is by exact forms supported near the
strata, hence the cohomology class is unchanged.
\end{theorem}

\begin{theorem}[Almost-complex normal-crossings model; Parker]
\label{thm:parker}
Let $(M,\Omega)$ be symplectic and let $D\subset M$ be an
$\Omega$-orthogonal symplectic simple normal crossings divisor.  Then there is
a neighbourhood $U$ of $D$ and an $\Omega$-compatible almost complex structure
$J_U$ on $U$ such that for every $p\in D$, if exactly the branches
$D_{i_1},\ldots,D_{i_r}$ pass through $p$, then there are smooth complex
coordinates
\[
(z_1,\ldots,z_n):U_p\longrightarrow \mathbb C^n
\]
on a neighbourhood $U_p\subset U$ of $p$ such that
\[
D\cap U_p=\{z_1\cdots z_r=0\},
\]
with $D_{i_a}\cap U_p=\{z_a=0\}$ for $a=1,\ldots,r$, and
\[
dz_a\circ J_U=i\,dz_a,\qquad a=1,\ldots,r .
\]

This is Parker's existence theorem for compatible almost-complex structures on
orthogonal symplectic normal-crossings divisors; see Parker
\cite[Section 4]{Parker}.  Parker's definition of an almost-complex
normal-crossings divisor is precisely the local-coordinate condition displayed
above.
\end{theorem}

We also recall the standard extension fact for compatible almost complex
structures.

\begin{lemma}[Extension of compatible almost complex structures]
\label{lem:extension}
Let $(M,\Omega)$ be symplectic, let $A\subset M$ be closed, and let $J_0$ be an
$\Omega$-compatible almost complex structure on a neighbourhood $U$ of $A$.
Then, after possibly shrinking $U$, there exists a global $\Omega$-compatible
almost complex structure $J$ on $M$ such that $J=J_0$ near $A$.
\end{lemma}

\begin{proof}
Choose open sets $A\subset V\Subset U$.  Let
\[
g_0(v,w)=\Omega(v,J_0w)
\]
on $U$.  This is a Riemannian metric on $U$.  Choose any global Riemannian
metric $g_1$ on $M$ and a cutoff function $\rho$ with $\rho\equiv 1$ on $V$ and
$\operatorname{supp}\rho\subset U$.  Then
\[
g=\rho g_0+(1-\rho)g_1
\]
is a global Riemannian metric equal to $g_0$ on $V$.

Define the $g$-skew-adjoint bundle endomorphism $A_g$ by
\[
\Omega(v,w)=g(A_gv,w).
\]
Then $-A_g^2$ is positive definite.  Let
\[
P=(-A_g^2)^{1/2}
\]
be its positive square root and set
\[
J=P^{-1}A_g .
\]
This is the standard polar-decomposition construction of the $\Omega$-compatible
almost complex structure associated to $g$; see for example McDuff--Salamon
\cite[Proposition 2.50]{MS}.  On $V$ we have $g=g_0$, hence $A_g=J_0$ and
therefore $J=J_0$.  Thus $J$ is global, $\Omega$-compatible, and agrees with
$J_0$ near $A$.
\end{proof}

\section*{Main theorem}

\begin{theorem}
Let $(M,\omega)$ be closed symplectic and let $D\subset M$ be a symplectic
simple normal crossings divisor.  After replacing $\omega$ by a symplectic form
connected to it through a cohomologous path along which $D$ remains a
symplectic simple normal crossings divisor, there exists an $\omega$-tame
almost complex structure $J$ such that, locally on $M$, there are smooth complex
coordinates $(z_1,\ldots,z_n)$ in which, for some $0\le r\le n$, the functions
$z_1,\ldots,z_r$ are $J$-holomorphic and
\[
D=\{z_1\cdots z_r=0\}.
\]
\end{theorem}

\begin{proof}
Apply Theorem \ref{thm:orthogonalization}.  We obtain a cohomologous deformation
$\omega_t$ of $\omega$ such that $D$ remains symplectic normal crossings for all
$t$ and is $\omega_1$-orthogonal.

Now apply Theorem \ref{thm:parker} to $(M,\omega_1,D)$.  This gives a
neighbourhood $U$ of $D$ and an $\omega_1$-compatible almost complex structure
$J_U$ on $U$ with the required local holomorphic normal-crossings coordinates
near every point of $D$.

By Lemma \ref{lem:extension}, $J_U$ extends to a global $\omega_1$-compatible,
hence $\omega_1$-tame, almost complex structure $J$ on $M$, unchanged near
$D$.  Therefore the local holomorphic coordinate description near $D$ remains
valid for the global $J$.

At a point $p\notin D$, no branch of $D$ passes through $p$, so we take
$r=0$ and choose any smooth complex coordinate chart on the underlying
$2n$-manifold.  Thus the displayed local model holds everywhere on $M$.

Finally, since the question permits replacing $\omega$ by the cohomologous
symplectic form $\omega_1$, the global almost complex structure $J$ is tame for
the allowed symplectic form.  This proves the claim.
\end{proof}

\begin{thebibliography}{9}

\bibitem{FTZ}
E. Farajzadeh Tehrani and A. Zinger,
\emph{On symplectic normal crossings divisors},
J. Symplectic Geom. \textbf{17} (2019), no. 6, 1495--1583.

\bibitem{McLean}
M. McLean,
\emph{The growth rate of symplectic homology and affine varieties},
Geom. Funct. Anal. \textbf{22} (2012), no. 2, 369--442.

\bibitem{MS}
D. McDuff and D. Salamon,
\emph{Introduction to Symplectic Topology},
3rd ed., Oxford University Press, 2017.

\bibitem{Parker}
B. Parker,
\emph{Exploded manifolds},
Adv. Math. \textbf{229} (2012), no. 6, 3256--3319.

\end{thebibliography}

\end{document}
```